% =====================================================================
% P5 merge snippet — FM01: abstract + keywords (REPLACEMENT CANDIDATE)
% Source: manuscript_v2 frontmatter/abstract.tex, re-scoped for the merged paper.
% The final paper already has an abstract (~250 words) and 11 keywords, so this is not a
% missing piece: it is a drop-in replacement for the case where the merged scope is adopted,
% because it names the three things the merge adds — the exact-vs-reduced energy flux, the
% five-axis taxonomy, and the verification register. Word count: 258 (script-checked).
% Word count verified by tools/check_snippets.py.
% Suggested placement: replaces the abstract in main.tex L70--75, keeping the graphical
% abstract above it unchanged.
% =====================================================================

\begin{abstract}
Wave propagation in solids that conduct heat, deform, and carry further fields such as electric
polarization, magnetic induction, diffusing species or microstructural degrees of freedom is governed by a
family of coupled models that is usually presented one member at a time. This review organises that family
around a single heat-flux template with three lags and one memory term, shows how the classical,
Lord--Shulman, Green--Lindsay, Green--Naghdi, dual-phase-lag, three-phase-lag, Moore--Gibson--Thompson,
fractional and memory-dependent theories arise (or fail to arise) from it, and records which limits are
regular and which singular. Every quantitative figure is generated by an accompanying script, and the
identities it rests on are checked symbolically or numerically. For the coupled problem in a linear
homogeneous medium we derive and verify a dispersion polynomial in
$K=k^2$, state where delays, fractional powers and memory kernels make the relation transcendental, and
use the polynomial to compare phase speed, attenuation and thermoelastic damping across theories. Two
objects that the literature often conflates are kept apart: the availability flux, which carries a
dissipation identity and is the right flux for interface energy partition, and the exact Eulerian
total-energy flux, whose first-order reduction is the flux usually quoted. A five-axis taxonomy then
places any candidate model, hybrids included, on one set of coordinates and states precisely when two studies are
comparable at all. Transient problems, surface, Lamb and bar waves, interfaces and layered media, size effects and
multiphysics couplings are treated in turn, and the review closes with a comparative synthesis, a compact pitfalls
checklist, a decision agenda, and a per-equation verification register that marks each registered equation
as derived, conditional, source-checked, unverified or open.
\medskip

\noindent\textbf{Keywords:} generalized thermoelasticity; second sound; dual-phase-lag;
Moore--Gibson--Thompson; fractional heat conduction; memory-dependent derivative; nonlocal elasticity;
piezoelectricity; Rayleigh waves; thermoelastic damping; interface energy partition; energy flux;
model taxonomy; verification register.
\end{abstract}
